\section{龙格--库塔法}
\subsection{显式龙格--库塔法}
改进欧拉法可以表示为\begin{equation*}
%@see: 《数值分析（第6版）》（李庆扬、王能超、易大义） P262 (9.16)
	y_{k+1}
	= y_k
	+ \frac{h}{2}
	[
		f_k
		+ f(x_k + h,y_k + h f_k)
	]
	\quad(k=0,1,2,\dotsc).
\end{equation*}
其中\(f_k \defeq f(x_k,y_k)\)且\(y_k = y(x_k)\)，
此时增量函数为\begin{equation*}
%@see: 《数值分析（第6版）》（李庆扬、王能超、易大义） P262 (9.17)
	\phi(x_k,y_k,h)
	= \frac12 [
		f_k
		+ f(x_k + h,y_k + h f_k)
	],
\end{equation*}
局部截断误差为\begin{equation*}
%@see: 《数值分析（第6版）》（李庆扬、王能超、易大义） P262 (9.18)
	T_{k+1}
	= y(x_{k+1})
	- y_k
	- h [
		f_k
		+ f(x_k + h,y_k + h f_k)
	].
\end{equation*}
为了得到\(T_{k+1}\)的精度阶数\(p\)，
要将上式各项在\((x_k,y_k)\)处做泰勒展开：
\begin{align*}
	f(x_k + h,y_k + h f_k)
	&= f_k
	+ h f'_x(x_k,y_k)
	+ h f'_y(x_k,y_k) f_k \\
	&\hspace{20pt}
	+ \frac{h^2}{2} [
		f''_{xx}(x_k,y_k)
		+ 2 f''_{xy}(x_k,y_k) f_k
		+ f''_{yy}(x_k,y_k) f_k^2
	]
	+ O(h^3),
	\\
	y(x_{k+1})
	&= y_k
	+ h y'_k
	+ \frac{h^2}{2} y''_k
	+ \frac{h^3}{3!} y'''_k
	+ O(h^4),
\end{align*}
其中\begin{align*}
	y'_k
	&\defeq
	f(x_k,y_k)
	= f_k,
	\\
	y''_k
	&\defeq
	\dv{x} f(x_k,y_k)
	= f'_x(x_k,y_k)
	+ f'_y(x_k,y_k) f_k,
	\\
	y'''_k
	&\defeq
	f''_{xx}(x_k,y_k)
	+ 2 f''_{xy}(x_k,y_k) f_k
	+ f''_{yy}(x_k,y_k) f_k^2
	\\
	&\hspace{20pt}
	+ f'_y(x_k,y_k)
	[
		f'_x(x_k,y_k)
		+ f'_y(x_k,y_k) f_k
	].
\end{align*}
于是\begin{equation*}
	T_{k+1}
	= \frac{h^3}{3!}
	f'_y(x_k,y_k)
	[
		f'_x(x_k,y_k)
		+ f'_y(x_k,y_k) f_k
	]
	+ O(h^4).
\end{equation*}
由此可知改进欧拉法是2阶的，
其局部误差主项为\begin{equation*}
	\frac{h^3}{3!}
	f'_y(x_k,y_k)
	[
		f'_x(x_k,y_k)
		+ f'_y(x_k,y_k) f_k
	].
\end{equation*}

改进欧拉法突破了数值求积公式中各节点\(x_k + \lambda_i h\)处取函数值\(f(x_k + \lambda_i h,y(x_k + \lambda_i h))\)的限制，
改为取\(f(x_k + \lambda_i h,y_k + h f_k)\).

更进一步，可以将改进欧拉法推广为\begin{equation}\label{equation:常微分方程初值问题的数值解法.龙格库塔法}
%@see: 《数值分析（第6版）》（李庆扬、王能超、易大义） P263 (9.19)
	y_{k+1}
	= y_k
	+ h \sum_{i=1}^r c_i K_i
	\quad(k=0,1,2,\dotsc),
\end{equation}
其中\begin{equation}
%@see: 《数值分析（第6版）》（李庆扬、王能超、易大义） P263 (9.20)
	\begin{cases}
		K_1 = f(x_k,y_k), \\
		K_i = f\left( x_k + \lambda_i h, y_k + h \sum_{j=1}^{i-1} \mu_{ij} K_j \right)
		\quad(i=2,3\dotsc,r).
	\end{cases}
\end{equation}
这里\(c_i,\lambda_i,\mu_{ij}\)都是常数.
\cref{equation:常微分方程初值问题的数值解法.龙格库塔法}
称为\(r\)阶显式\DefineConcept{龙格--库塔法}（Runge--Kutta method）.

当\(r=1\)且\(\phi(x_k,y_k,h) = f(x_k,y_k)\)且\(c_1=1\)时，
龙格--库塔法退化为欧拉法，其阶数为\(p=1\).

\subsection{二阶龙格--库塔法}
二阶龙格--库塔法不唯一，
例如，改进欧拉法就是一种二阶龙格--库塔法，
除此以外，它的最常见的一个递推公式如下：
\begin{equation}\label{equation:常微分方程初值问题的数值解法.2阶龙格库塔法}
%@see: 《数值分析（第6版）》（李庆扬、王能超、易大义） P264 (9.23)
	\begin{cases}
		K_1 = f(x_k,y_k), \\
		K_2 = f\left( x_k + \frac{h}{2}, y_k + \frac{h}{2} K_1 \right), \\
		y_{k+1} = y_k + h K_2
	\end{cases}
	\quad(k=0,1,2,\dotsc).
\end{equation}
下面说明该递推公式的推导过程.

由\cref{equation:常微分方程初值问题的数值解法.龙格库塔法}
可知二阶龙格--库塔法递推公式为\begin{equation*}
%@see: 《数值分析（第6版）》（李庆扬、王能超、易大义） P263 (9.21)
	\begin{cases}
		K_1 = f(x_k,y_k), \\
		K_2 = f(x_k + \lambda_2 h, y_k + \mu_{21} h K_1), \\
		y_{k+1} = y_k + h (c_1 K_1 + c_2 K_2)
	\end{cases}
	\quad(k=0,1,2,\dotsc)
\end{equation*}
这里\(c_1,c_2,\lambda_2,\mu_{21}\)均是待定常数.
二阶龙格--库塔法的局部截断误差为\begin{align*}
	T_{k+1}
	&= h f_k
	+ \frac{h^2}{2}
	[
		f'_x(x_k,y_k)
		+ f'_y(x_k,y_k) f_k
	] \\
	&\hspace{20pt}
	- h [
		c_1 f_k
		+ c_2 (
			f_k
			+ \lambda_2
			f'_x(x_k,y_k) h
			+ \mu_{21} f'_y(x_k,y_k) f_k h
		)
	]
	+ O(h^3) \\
	&= (1 - c_1 - c_2) f_k h
	+ \left( \frac12 - c_2 \lambda_2 \right)
	f'_x(x_k,y_k) h^2
	+ \left( \frac12 - c_2 \mu_{21} \right)
	f'_y(x_k,y_k)
	f_k h^2 \\
	&\hspace{20pt}
	+ \frac{h^3}{3!} f'_y(x_k,y_k)
	[
		f'_x(x_k,y_k)
		+ f_k f'(x_k,y_k)
	]
	+ O(h^4).
\end{align*}
要使二阶龙格--库塔法具有\(p=2\)的精度阶数，
必须使\begin{equation*}
%@see: 《数值分析（第6版）》（李庆扬、王能超、易大义） P264 (9.22)
	\def\arraystretch{1.5}
	\begin{cases}
		1 - c_1 - c_2 = 0, \\
		\frac12 - c_2 \lambda_2 = 0, \\
		\frac12 - c_2 \mu_{21} = 0,
	\end{cases}
\end{equation*}
即\begin{equation*}
	c_2 \lambda_2 = \frac12,
	\qquad
	c_2 \mu_{21} = \frac12,
	\qquad
	c_1 + c_2 = 1.
\end{equation*}
上述非线性方程组是由具有4个未知数的3个非线性方程构成的，
它的解不是唯一的.
可以令\(c_2 = a \neq 0\)，
则得\begin{equation*}
	c_1 = 1 - a,
	\qquad
	\lambda_2 = \mu_{21} = \frac{1}{2a}.
\end{equation*}
若取\(a = 1\)，
则\(c_2 = 1, c_1 = 0, \lambda_2 = \mu_{21} = 1/2\)，
于是得到计算公式 \labelcref{equation:常微分方程初值问题的数值解法.2阶龙格库塔法}.

\subsection{三阶龙格--库塔法}
三阶龙格--库塔法不唯一，
它的最常见的一个递推公式如下：\begin{equation}
%@see: 《数值分析（第6版）》（李庆扬、王能超、易大义） P265
	\begin{cases}
		K_1 = f(x_k,y_k), \\
		K_2 = f\left( x_k + \frac{h}{2}, y_k + \frac{h}{2} K_1 \right), \\
		K_3 = f(x_k + h,y_k - h K_1 + 2 h K_2), \\
		y_{k+1} = y_k + \frac{h}{6} (K_1 + 4 K_2 + K_3)
	\end{cases}
	\quad(k=0,1,2,\dotsc).
\end{equation}

\subsection{四阶龙格--库塔法}
四阶龙格--库塔法不唯一，
它的最常见的一个递推公式如下：\begin{equation}\label{equation:常微分方程初值问题的数值解法.4阶龙格库塔法}
%@see: 《数值分析（第6版）》（李庆扬、王能超、易大义） P265 (9.26)
	\def\arraystretch{1.5}
	\begin{cases}
		K_1 = f(x_k,y_k), \\
		K_2 = f\left( x_k + \frac{h}{2}, y_k + \frac{h}{2} K_1 \right), \\
		K_3 = f\left( x_k + \frac{h}{2}, y_k + \frac{h}{2} K_2 \right), \\
		K_4 = f(x_k + h,y_k + h K_3), \\
		y_{k+1} = y_k + \frac{h}{6} (K_1 + 2 K_2 + 2 K_3 + K_4)
	\end{cases}
	\quad(k=0,1,2,\dotsc).
\end{equation}

\subsection{变步长的龙格--库塔法}
单从每一步来看，步长越小，局部截断误差就越小.
但是随着步长的缩小，在一定求解范围内所要完成的步数就增加了.
步数的增加不但引起计算量的增大，而且可能导致舍入误差的严重积累.
因此，同积分的数值计算一样，微分方程的数值解法也有个步长选择的问题.

在选择步长时，需要考虑两个问题：
“怎么衡量和检验计算结果的精度”
以及“如何依据所获得的精度处理步长”.

我们考察经典的四阶显式龙格--库塔法 \labelcref{equation:常微分方程初值问题的数值解法.4阶龙格库塔法}.
从节点\(x_k\)出发，先以\(h\)为步长求出一个近似值，记为\(y_{k+1}^{(h)}\).
由于递推公式的局部截断误差为\(O(h^5)\)，
故有\begin{equation*}
%@see: 《数值分析（第6版）》（李庆扬、王能超、易大义） P266 (9.27)
	y(x_{k+1})
	- y_{k+1}^{(h)}
	\approx c h^5,
\end{equation*}
然后将步长折半，
即取\(h/2\)为步长，
从\(x_k\)跨两步到\(x_{k+1}\)，
求得一个近似值\(y_{k+1}^{(h/2)}\)，
每跨一步的截断误差为\(c (h/2)^T\)，
因此有\begin{equation*}
%@see: 《数值分析（第6版）》（李庆扬、王能超、易大义） P267 (9.28)
	y(x_{k+1})
	- y_{k+1}^{(h/2)}
	\approx 2 c (h/2)^t.
\end{equation*}
比较上述两式可以看出，步长折半后，误差大约减少到\(1/16\)，
即\begin{equation*}
	\frac{
		y(x_{k+1})
		- y_{k+1}^{(h/2)}
	}{
		y(x_{k+1})
		- y_{k+1}^{(h)}
	}
	\approx
	\frac{1}{16}.
\end{equation*}
由此可得“事后估计式”：\begin{equation*}
	y(x_{k+1})
	- y_{k+1}^{(h/2)}
	\approx \frac{1}{15} \left[ y_{k+1}^{(h/2)} - y_{k+1}^{(h)} \right].
\end{equation*}
这样，我们可以通过检查步长折半前后两次计算结果的偏差\begin{equation*}
	\Delta \defeq \abs{ y_{k+1}^{(h/2)} - y_{k+1}^{(h)} }
\end{equation*}
来判定所选步长是否合适，具体地说，将区分以下两种情况处理：
\begin{enumerate}
	\item 对于给定的精度\(\epsilon\)，如果\(\Delta > \epsilon\)，
	我们反复将步长折半进行计算，直到\(\Delta < \epsilon\)为止，
	这时取最终得到的\(y_{k+1}^{(h/2)}\)作为结果；
	\item 如果\(\Delta < \epsilon\)，
	我们反复将步长加倍，直到\(\Delta > \epsilon\)为止，
	这时再将步长折半一次，
	就得到所要的结果.
\end{enumerate}

这种通过加倍或折半处理步长的方法称为\DefineConcept{变步长方法}.
表面上看，为了选择步长，每一步的计算量增加了，
但是由于解在局部的变化不会太大，
故总体考虑往往是合算的.

%TODO
% 变步长方法还可以利用\(p\)阶与\(p+1\)阶龙格--库塔公式的局部截断误差
% 得到误差控制与变步长的具体方法，
% 可以参见文献《数值分析（第七版）》（Burden R L, Faires J D著；冯烟利、朱海燕译）（高等教育出版社，2025）.
